% Persona vs.\ Model Variance in LLM-Elicited Expert Panels
% This section reports the experimental design and results for the
% persona-diversity and cross-model variance analyses.

\subsection{Overview and Research Questions}
\label{sec:persona-experiment-design}

LLM-simulated Delphi panels assume that assigning distinct expert personas to the same language model produces meaningfully diverse probability estimates. We test this assumption directly and contrast it with a positive control: whether \emph{model identity} generates significant variance.

\textbf{Research questions:}
\begin{enumerate}
  \item Does persona assignment systematically shift probability estimates, or are differences indistinguishable from sampling noise?
  \item Does model identity generate significant variance, validating our statistical methodology?
\end{enumerate}

\paragraph{Experimental setup.}
\textbf{Scenario:} OC3 ransomware targeting a large enterprise (\$250M--\$1B revenue, MIL2/MIL3 cybersecurity maturity). LLMs estimate the probability of attacker success at each MITRE ATT\&CK step, conditional on completing the Imaginairy benchmark (BountyBench).

\textbf{Attack steps} (spanning different baseline probabilities):
\begin{itemize}
  \item \textbf{TA0002 -- Execution} (50\% baseline)
  \item \textbf{TA0007 -- Discovery} (85\% baseline)
  \item \textbf{T1657 -- Financial Theft} (30\% baseline)
\end{itemize}

\textbf{Personas:} Ten cybersecurity experts with distinct backgrounds---Defensive Security Specialist, Malware Reverse Engineer, AI/ML Security Researcher, Threat Intelligence Analyst, Security Software Developer, Incident Response Specialist, Academic Security Researcher, Red Team Operator, CTF Competition Expert, and Security Compliance Officer---each defined via system prompts specifying background, analytical style, and domain biases.

\textbf{Models:} Claude Sonnet~4.5, GPT-4o, and Gemini~2.5~Pro. All use temperature $T=1.0$ with a single Delphi round (no deliberation), ensuring statistical independence of estimates.


\subsection{Experiment 1: Does Persona Assignment Matter?}
\label{sec:persona-variance}

\paragraph{Design.}
For each model--step combination: $10 \times 10 = 100$ probability estimates (10 personas $\times$ 10 independent runs). Each run executes all personas concurrently with no shared context---a balanced one-way ANOVA with 10 groups and 10 observations per group.

\paragraph{Statistical approach.}
We decompose total variance into between-persona variance ($\sigma^2_B$) and within-persona variance ($\sigma^2_W$, LLM sampling noise). The null hypothesis $H_0$ states all persona means are equal ($\sigma^2_B = 0$).

\textbf{Test statistics:}
\begin{enumerate}
  \item \textbf{ANOVA $F$-test:} Tests whether between-persona variance exceeds within-persona variance.
  \item \textbf{Kruskal--Wallis $H$:} Non-parametric alternative, robust to non-normality.
  \item \textbf{ICC(1,1):} Proportion of variance explained by persona identity (0 = none, 1 = all).\footnote{$\mathrm{ICC} = (\mathrm{MS}_B - \mathrm{MS}_W) / (\mathrm{MS}_B + (n{-}1)\,\mathrm{MS}_W)$}
\end{enumerate}

\paragraph{Primary result.}
\textbf{Finding: Persona assignment does not generate meaningful variance.}

Table~\ref{tab:persona-primary} shows results for Claude Sonnet~4.5, TA0002 (50\% baseline). The 10 persona means span only 3.2 percentage points (pp), from 0.512 to 0.544. Within-persona variability (SD = 0.025) exceeds between-persona variability (SD = 0.010).

The ANOVA yields $F = 1.61$, $p = 0.12$---\emph{not significant}. ICC = 0.058 means persona identity explains only 5.8\% of variance; 94.2\% is sampling noise. With 100 observations, this design has $>$95\% power to detect even a 2pp spread; the observed 1pp spread is half this threshold. \textbf{Personas are indistinguishable from noise.}

\begin{table}[t]
\centering
\caption{Primary persona experiment (Claude Sonnet~4.5, TA0002, 50\% baseline). Ten personas $\times$ 10 runs = 100 estimates. Persona means span only 3.2pp.}
\label{tab:persona-primary}
\small
\begin{tabular}{lc}
\toprule
\textbf{Persona} & \textbf{Mean} \\
\midrule
AI/ML Security Researcher      & 0.512 \\
Malware Reverse Engineer       & 0.517 \\
Incident Response Specialist   & 0.521 \\
Defensive Security Specialist  & 0.522 \\
Red Team Operator              & 0.524 \\
Threat Intelligence Analyst    & 0.524 \\
CTF Competition Expert         & 0.528 \\
Security Software Developer    & 0.534 \\
Security Compliance Officer    & 0.539 \\
Academic Security Researcher   & 0.544 \\
\midrule
\multicolumn{2}{l}{\textbf{Variance:} Within-persona SD = 0.025 (noise), Between-persona SD = 0.010 (signal)} \\
\midrule
\multicolumn{2}{l}{\textbf{Tests:} $F = 1.61$, $p = 0.12$ (not sig.); $H = 11.27$, $p = 0.26$; ICC = 0.058} \\
\bottomrule
\end{tabular}
\end{table}

\paragraph{Robustness check.}
We replicate the $10 \times 10$ design across 2 models $\times$ 3 attack steps = 6 experiments (Table~\ref{tab:persona-robustness}).

\textbf{Result: 5 of 6 experiments show no significant persona effect.} The single nominally significant result (Claude, TA0007, $p = 0.01$) has negligible practical effect: persona means span only 2.1pp (0.854--0.875), ICC = 0.144 (85.6\% variance is noise), does not replicate on GPT-4o ($p = 0.15$), and fails Bonferroni correction ($\alpha_{\text{corrected}} = 0.008$).

\begin{table}[t]
\centering
\caption{Persona ANOVA robustness (6 experiments, each 10 personas $\times$ 10 runs). $^*$Does not survive Bonferroni correction ($\alpha = 0.008$).}
\label{tab:persona-robustness}
\small
\begin{tabular}{llccc}
\toprule
\textbf{Model} & \textbf{Step (Baseline)} & $\boldsymbol{F}$ & $\boldsymbol{p}$ & \textbf{ICC} \\
\midrule
Claude Sonnet 4.5 & TA0002 (50\%) & 1.61 & 0.12  & 0.058 \\
GPT-4o            & TA0002 (50\%) & 0.77 & 0.65  & 0.000 \\
\midrule
Claude Sonnet 4.5 & TA0007 (85\%) & 2.51 & 0.01$^*$ & 0.144 \\
GPT-4o            & TA0007 (85\%) & 1.54 & 0.15  & 0.051 \\
\midrule
Claude Sonnet 4.5 & T1657\phantom{0} (30\%) & 0.50 & 0.87  & 0.000 \\
GPT-4o            & T1657\phantom{0} (30\%) & 1.11 & 0.36  & 0.011 \\
\bottomrule
\end{tabular}
\end{table}


\subsection{Experiment 2: Does Model Identity Matter? (Positive Control)}
\label{sec:cross-model-variance}

\paragraph{Motivation.}
If our ANOVA framework correctly identifies that personas do \emph{not} matter, can it detect factors that \emph{do}? We test model identity as a positive control.

\paragraph{Design.}
For each attack step, all three models produce $10 \times 10 = 100$ estimates. Since persona $\approx$ noise (Experiment 1), we treat the 100 per-model observations as independent draws. One-way ANOVA with model as grouping factor (3 groups, $\sim$300 observations per step).

\paragraph{Results.}
\textbf{Finding: Model identity generates highly significant variance.}

Table~\ref{tab:cross-model} shows all three experiments yield $p < 0.0001$ (Bonferroni-corrected $\alpha = 0.017$). Model identity explains 28--65\% of variance (ICC) vs.\ 0--14\% for personas. Model means span 2.5--12.2pp, far exceeding the $\leq$3.2pp persona spread.

\begin{table}[t]
\centering
\caption{Cross-model ANOVA (3 models $\times$ $\sim$100 estimates = $\sim$300 observations per step). All highly significant.}
\label{tab:cross-model}
\small
\begin{tabular}{lcccccc}
\toprule
\textbf{Step} & \textbf{Claude} & \textbf{GPT-4o} & \textbf{Gemini} & $\boldsymbol{F}$ & $\boldsymbol{p}$ & \textbf{ICC} \\
\midrule
TA0002 (50\%) & 0.527 & 0.648 & 0.556 & 187.79 & ${<}0.0001$ & 0.651 \\
TA0007 (85\%) & 0.863 & 0.888 & 0.885 & \phantom{0}38.29 & ${<}0.0001$ & 0.279 \\
T1657\phantom{0} (30\%) & 0.323 & 0.381 & 0.383 & 155.35 & ${<}0.0001$ & 0.607 \\
\midrule
\multicolumn{7}{l}{Model means span 2.5--12.2pp vs.\ $\leq$3.2pp for personas.} \\
\bottomrule
\end{tabular}
\end{table}

\paragraph{Direct comparison: Persona vs.\ Model.}
Table~\ref{tab:persona-vs-model} contrasts persona and model effects on the same attack steps. \textbf{Model $F$-statistics are 15--311$\times$ larger than persona $F$-statistics.} Model ICC (0.28--0.65) vastly exceeds persona ICC (0.00--0.14). The same ANOVA framework that returns null for personas produces overwhelming evidence for model effects---validating the methodology.

\begin{table}[t]
\centering
\caption{Persona vs.\ model: Same ANOVA, opposite conclusions. Model effects are 15--311$\times$ stronger.}
\label{tab:persona-vs-model}
\small
\begin{tabular}{lcccccc}
\toprule
 & \multicolumn{2}{c}{$\boldsymbol{F}$-statistic} & \multicolumn{2}{c}{\textbf{ICC (\% variance)}} & \multicolumn{2}{c}{\textbf{Effect}} \\
\cmidrule(lr){2-3} \cmidrule(lr){4-5} \cmidrule(lr){6-7}
\textbf{Step} & Persona & Model & Persona & Model & Persona & Model \\
\midrule
TA0002 (50\%) & 1.61 & 187.79 & 5.8\% & 65.1\% & Null & \textbf{Strong} \\
TA0007 (85\%) & 2.51 & \phantom{0}38.29 & 14.4\% & 27.9\% & Null & \textbf{Strong} \\
T1657\phantom{0} (30\%) & 0.50 & 155.35 & 0.0\% & 60.7\% & Null & \textbf{Strong} \\
\bottomrule
\end{tabular}
\end{table}


\subsection{Summary}
\label{sec:persona-summary}

\textbf{Key findings:}
\begin{enumerate}
  \item \textbf{Persona assignment does not matter.} Across 6 experiments (2 models $\times$ 3 baselines, 600 estimates), persona explains $\leq$14\% of variance. Five of six show no significant effect; the one nominally significant result fails Bonferroni correction and does not replicate.

  \item \textbf{Model identity matters.} Across 3 experiments (3 models $\times$ 3 baselines, $\sim$890 estimates), model choice explains 28--65\% of variance (all $p < 0.0001$). Model $F$-statistics are 15--311$\times$ larger than persona $F$-statistics.

  \item \textbf{Methodology validated.} The same ANOVA framework correctly identifies persona as a null effect and model as a strong effect, confirming the design is well-powered and the persona null is a true negative.
\end{enumerate}

\textbf{Implication:} For LLM-elicited expert panels, \emph{which LLM to use} matters far more than \emph{which persona to assign}. Multi-model ensembles may be necessary to capture this uncertainty.
